% ============================================================
\section{Conclusion}
\label{sec:conclusion}
% ============================================================

We investigated whether synthetic hard negatives improve contrastive
anomaly detection for IoT IDS, using a five-condition ablation with
protocol-constrained stealth-attack synthesis and benign-calibrated thresholds.
Three key findings emerge.

First, the NLL$+$SupCon training recipe is \emph{backbone-agnostic}:
it improves detection for both Moirai (C, AUC$=$0.563) and a from-scratch
CNN (AUC$=$0.597) at 200/5, and generalises cross-domain to N-BaIoT
(AUC$=$\nbaiotAUC{}).
Second, the C$>$D gap (Gaussian-noise vs.\ hard negatives) collapses
to near-zero at 5$\times$ scale with NLL-only early stopping---a
small-scale training artefact: the advantage of easy negatives diminishes
with sufficient data.
Third, \textbf{catastrophic forgetting} is the dominant failure mode at
extended scale: freezing the encoder recovers AUC
(frozen B$=$0.606, C$=$0.582, D$=$0.573 vs.\ unfrozen 0.506--0.521),
confirming that fine-tuning destroys the pre-trained features that
provide discriminative power.

The practical implication: freeze the encoder when fine-tuning foundation
models at scale; use CNN at small scale; recalibrate thresholds on
deployment-specific benign traffic.
We release all code, datasets, and calibration scripts.
